% ─────────────────────────────────────────────────────────────────────────────
% front/abstract.tex — Three abstracts: French, English, Arabic
% ─────────────────────────────────────────────────────────────────────────────

% ── Résumé (French) ───────────────────────────────────────────────────────────
\phantomsection
\addcontentsline{toc}{chapter}{Résumé}
\chapter*{Résumé}

Le présent mémoire décrit la conception et le développement de \textit{Mobile Match Algérie}, une application web permettant aux consommateurs algériens de comparer les offres mobiles des trois principaux opérateurs~: Mobilis, Djezzy et Ooredoo. Face à la multiplicité et à la complexité croissante des forfaits proposés sur le marché national, les utilisateurs manquent d'un outil neutre et centralisé pour prendre des décisions éclairées. Ce travail répond à ce besoin en proposant une plateforme regroupant, en temps réel, l'ensemble des offres disponibles grâce à un module de collecte automatisée~; en offrant des outils de filtrage, de comparaison côte à côte et de visualisation graphique~; et en personnalisant les recommandations selon le profil budgétaire et d'usage de chaque utilisateur.

L'application est développée avec le framework \textit{Next.js~16}, en utilisant \textit{Prisma~ORM} avec une base de données \textit{SQLite}, et déployée en tant qu'application full-stack. Un moteur de recommandation basé sur un algorithme de scoring multi-critères sélectionne les trois offres les mieux adaptées au profil de l'utilisateur. L'ensemble des fonctionnalités définies dans le cahier des charges a été implémenté et validé fonctionnellement.

\bigskip
\noindent\textbf{Mots-clés~:} comparaison d'offres mobiles, web scraping, système de recommandation, Next.js, Algérie.

\newpage

% ── Abstract (English) ───────────────────────────────────────────────────────
\phantomsection
\addcontentsline{toc}{chapter}{Abstract}
\chapter*{Abstract}

This thesis presents the design and development of \textit{Mobile Match Algeria}, a web application that enables Algerian consumers to compare mobile plans offered by the country's three main operators: Mobilis, Djezzy, and Ooredoo. The increasing number and complexity of available plans make it difficult for users to make informed choices without a neutral, centralised comparison tool. This work addresses that gap by delivering a platform that aggregates all current offers in real time through an automated scraping module, provides advanced filtering, side-by-side comparison, and graphical visualisation tools, and personalises plan recommendations according to each user's budget and usage profile.

The application is built with the \textit{Next.js~16} framework using \textit{Prisma~ORM} backed by \textit{SQLite}, and is deployed as a full-stack web service. A recommendation engine based on a multi-criteria scoring algorithm selects the three best-matched plans for any given user profile. All functional requirements defined in the specification have been implemented and verified.

\bigskip
\noindent\textbf{Keywords:} mobile plan comparison, web scraping, recommendation system, Next.js, Algeria.

\newpage

% ── ملخص (Arabic) ────────────────────────────────────────────────────────────
\clearpage
\thispagestyle{plain}
\phantomsection
\addcontentsline{toc}{chapter}{\texorpdfstring{{\arabicfont ملخص}}{Mulakhkhas}}

\begin{RTL}
{\arabicfont\onehalfspacing

\begin{center}
  {\bfseries\fontsize{14}{17}\selectfont ملخص}\\[8pt]
  {\large\bfseries Mobile Match Algérie}
\end{center}

\bigskip

يُقدِّم هذا المذكرة تصميم وتطوير تطبيق ويب يُتيح للمستهلكين الجزائريين مقارنة
عروض الهاتف المحمول التي تقدمها المشغلات الثلاثة الرئيسية: موبيليس وجيزي وأوريدو.
يُعالج هذا العمل غياب أداة مقارنة محايدة ومركزية من خلال منصة تجمع جميع العروض
الحالية في الوقت الفعلي، وتوفر أدوات تصفية متقدمة ومقارنة جانبية وتصوراً بيانياً،
مع تخصيص التوصيات وفق الميزانية وملف استخدام كل مستخدم.

\bigskip

بُني التطبيق بإطار العمل Next.js 16 مع Prisma ORM وقاعدة بيانات SQLite،
ونُشر بوصفه خدمة ويب متكاملة. يعتمد محرك التوصيات على خوارزمية تسجيل متعددة
المعايير تختار أفضل ثلاثة عروض تناسب ملف المستخدم. وقد استُكملت جميع المتطلبات
الوظيفية المحددة في دفتر الشروط وتم التحقق منها.

\bigskip
\noindent\textbf{الكلمات المفتاحية:}
مقارنة عروض الهاتف المحمول،
استخلاص البيانات من الويب،
نظام التوصيات،
Next.js،
الجزائر.

}
\end{RTL}

\newpage
